% =============================================================================
% §3 PROBLEM FORMULATION — page budget: 0.75 pages (paper_plan_v2 §9)
%
% Content sources:
%   * episodic decision process — plan §2.1: ANSWER = stop action; T_max;
%     forced-continuation vs RL rollout modes (the critical distinction);
%     x_t = inference-available features ONLY; q_t per domain; multi-dimensional
%     dollar cost c_t; pilot-normalized C-tilde; U_t = q_t − λ·m(tier)·c̃.
%   * Snell envelope — plan §2.2 / §10 Alg.1: V_T = U_T;
%     V_t = max(U_t, E[V_{t+1}|x_t]); Cont; Δ*; τ* = min{t: U_t ≥ Cont(x_t)}.
%   * classical citations (CITED, NOT CLAIMED — plan §3 contribution 3):
%     Hansen & Zilberstein AIJ'01 (anytime monitoring), Weitzman '79,
%     Longstaff–Schwartz regression, prophet inequalities.
%   * prophet-bias motivation for the label choice — plan §2.2: argmax of the
%     realized future upper-bounds every non-anticipating rule => stops late;
%     max-of-mean (Snell) replaces mean-of-max.
% =============================================================================
\section{Problem Formulation}
\label{sec:formulation}

\TODO{Episodic loop with stop action; two rollout modes (label-collection
forced-continuation vs RL termination) --- plan \S2.1 verbatim structure.}

\TODO{Stopping utility $U_t = q_t - \sum_{i\le t}\lambda\, m(\mathrm{tier}_i)\,
\tilde c_i$; tier-scaled marginal costing as discretized Lagrangian relaxation
(plan \S2.2).}

\TODO{Empirical Snell envelope by backward recursion with cross-sectional
regression (Longstaff--Schwartz); define $\mathrm{Cont}(x_t)$, $\Delta^*_t$,
$\tau^*$; the three precision notes (decision grid, $t{=}T$ edge, policy-relative
$\tau^*$) --- plan \S2.2/\S10.}

\TODO{Prophet-bias argument for why the naive argmax label stops late (plan
\S2.2) --- classical roots cited, not claimed.}
